\begin{frame}{Зачем?}
  \begin{enumerate}
    \item Целые числа представляются как последовательность нулей и единиц
      в двоичной системе счисления
    \item Длина последовательности фиксирована --- 8, 16, 32, 64, \ldots бита
    \item А значит ограничены наибольшее и наименьшее число \pause
    \item Можно брать более длинные последовательности, но это не
      всегда оптимально
  \end{enumerate}
\end{frame}

\begin{frame}{Как?}{Идея}
  \begin{enumerate}
    \item Можно изменять количество членов последовательности
    \item Можно использовать не только двоичную систему счисления
  \end{enumerate}
\end{frame}

\begin{frame}{Kак?}{Реализация}
  \begin{enumerate}
    \item Будем использовать последовательности произвольной длины в куче
    \item Cистема счисления не двоичная, а $2^n$-чная, где $n$ ---
      длина машинного слова
  \end{enumerate}

  $$(2^{32}-1,\ldots)+(1,\ldots)=(0,1,\ldots)$$\pause
  $$(2^{32}-1,\ldots)+(2^{32}-1,\ldots)+(0,3,\ldots)=(2^{32} - 2,4,\ldots)$$
  $$\underbrace{(2^{32} - 1)*(2^{32})^0}_{(2^{32}-1,\ldots)} +
  \underbrace{(2^{32} - 1)*(2^{32})^0}_{(2^{32}-1,\ldots)} +
  \underbrace{3*(2^{32})^1}_{(0,3,\ldots)}
  = 2^{33} - 2 + 3*2^{32}$$
\end{frame}

\begin{frame}{GMP}
  GMP (GNU Multiple Precision Arithmetic library) --- библиотека на
  языке C, реализующая такую модель.
  \lstinputlisting[language=C]{gmp.c}
\end{frame}

\begin{frame}{GMP}{Где используют?}
  \begin{enumerate}
    \item OpenSSL и libgcrypt используют GMP для реализации криптографии
    \item Mathematica (система компьютерной алгебры)
    \item Integer в GHC (компилятор языка Haskell),
      \texttt{java.math.BigInteger} и \texttt{java.math.BigDecimal} в
      GNU Classpath
  \end{enumerate}
\end{frame}

\begin{frame}
  \begin{center}
    {\usebeamercolor[fg]{structure} \Huge Демонстрация}
  \end{center}
  \vspace{0.5cm}
  \begin{center}
    {\large Код и слайды в репозитории ниже}
    {\large \url{https://github.com/kislball/gmp-report}}
  \end{center}
\end{frame}
